\section{Theorem of Dimensional Response Confluence}
\label{rm:ch:confluencia-general}

\subsection{Antecedent common and product fibered}

The volume results arise from heterogeneous readers applied to
the same enriched genealogical state. Let

\begin{equation}
 X=\bigl(X_{\APP},X_{\TRIT},X_{\TPK},
          \mathcal C,\mathcal B,\mathcal R,\mathcal M,\partial X\bigr)
 \label{rm:eq:confluence-state}
\end{equation}

the state that preserves cylinders, carry, route, memory, and boundary. The typed
readers are constructed on \(X\)

\begin{equation}
 P_i:X\longrightarrow Y_i,
 \qquad
 i\in I:=\{\alpha,{\rm act},{\rm ang},108,\square,{\rm obs},{\rm Per}\}.
 \label{rm:eq:confluence-objects}
\end{equation}

Their codomains represent, respectively, the dodecaphasic
coupling generator, the action pair, the bilayer angular operator, the clock
\(108\), the cubic screen, the internal observer, and the Perron mode.
For each reader, its graph is taken, provided with the projection in its canonical form

\begin{equation}
 \Gamma_i:=\{(x,y_i)\in X\times Y_i:P_i(x)=y_i\},
 \qquad
 \pi_i:\Gamma_i\longrightarrow X,
 \quad (x,y_i)\longmapsto x.
 \label{rm:eq:confluence-reader-graphs}
\end{equation}

The master object is the well-typed fiber product of those graphs:

\begin{equation}
 \boxed{
 \mathfrak M_X=
 \mathop{\times}_{X,\,i\in I}\Gamma_i
 =\left\{\bigl(x,(y_i)_{i\in I}\bigr):
 P_i(x)=y_i\ \text{for every }i\in I\right\}.}
 \label{rm:eq:master-fiber-product}
\end{equation}

The projections \(\pi_i\) materially fix the same antecedent. An
element of \(\mathfrak M_X\) is a family of compatible outputs
accompanied by the state that generates them.

\begin{definition}[Publication dimensional]
A dimensional publication is a composition
\begin{equation}
 \mathfrak M_X\xrightarrow{\;P\;}\mathcal L
 \xrightarrow{\;\operatorname{ev}_{\mathcal U}\;}\RR,
 \label{rm:eq:dimensional-publication}
\end{equation}
where \(\mathcal L\) is a typed dimensional line, and
\(\operatorname{ev}_{\mathcal U}\) is the final evaluation in a chart of
units. The map \(P\) first constructs the invariant and its genealogy.
\end{definition}

\subsection{Horizontal independence and vertical compatibility}

Two sibling readers share an antecedent and may satisfy cross-relations
without one factoring through the other. This distinction governs the
double sheets of the vacuum, the positive and oriented realizations, the double
Archimedean and exceptional projection, and the geometric,
informational, and flavor readings of the transition \(6\to8\).

\begin{theorem}[Confluence without factorization horizontal]
Let \(P_i:\mathfrak M_X\to Y_i\) and \(i=1,\ldots,r\) be publications
constructed on the same state. Suppose that each \(P_i\) preserves its
typed data and that the relations \(R_j(P_1,\ldots,P_r)=0\) are proved
after the readers have been constructed. Then the joint morphism
\begin{equation}
 P=(P_1,\ldots,P_r):\mathfrak M_X\longrightarrow
 Y_1\times\cdots\times Y_r
 \label{rm:eq:joint-publication}
\end{equation}
is vertically compatible with \(X\). A factorization
\(P_i=F_{ij}\circ P_j\) additionally requires \(P_j\) to separate all the
fibers distinguished by \(P_i\). The mere existence of \(R_j=0\) does not imply
that separation.
\end{theorem}

\begin{proof}
Vertical compatibility follows from the universal property of the
fiber product of the graphs \(\Gamma_i\): all components have
the same image in \(X\).
If \(P_i=F_{ij}\circ P_j\), then
\(P_j(x)=P_j(x')\) forces \(P_i(x)=P_i(x')\). Hence, any pair
\(x,x'\) that shares the reading \(j\) and differs in the reading \(i\) refutes
the factorization. A scalar relation \(R_j=0\) only restricts the joint
image; it does not guarantee the inclusion of kernels required for factorization.
\end{proof}

The theorem explains why the product of sheets can fix the causal cone while the quotient preserves orientation, and why the total area can coexist with a modular Hamiltonian that distinguishes the origin \(90/120\).

\subsection{1st confluence: coupling, vacuum and the electroweak sector}

Let

\begin{equation}
 D=D_A=\det T,
 \qquad
 Q(D)=-\frac9{25}\log D=4\pi\alpha_{\HMT}.
 \label{rm:eq:confluence-q-d}
\end{equation}

The memory-oriented is expressed through

\begin{equation}
 \rho=e^{2C^*_{\rm rad}},
 \qquad
 q_-=\sqrt{D\rho},
 \qquad
 q_+=\sqrt{D/\rho}.
 \label{rm:eq:confluence-qpm}
\end{equation}

The functional calculus

\begin{equation}
 F(z)=\frac{1-z^{90}}{1-z^{120}},
 \qquad
 r_\pm=F(q_\pm)
 \label{rm:eq:confluence-F}
\end{equation}

yields

\begin{equation}
 \widehat\mu=r_-^2,
 \quad
 \widehat\varepsilon=r_+^2,
 \quad
 \widehat Z=\frac{r_-}{r_+},
 \quad
 \widehat c=(r_-r_+)^{-1}.
 \label{rm:eq:confluence-vacuum}
\end{equation}

At the tree-level electroweak interface,

\begin{equation}
 g^2=\frac{Q(D)}{1-D},
 \qquad
 g'^2=\frac{Q(D)}{D},
 \label{rm:eq:confluence-gauge-couplings}
\end{equation}

from which

\begin{equation}
 \boxed{
 (1-D)g^2=Dg'^2=Q(D),
 \qquad
 \frac1{g^2}+\frac1{g'^2}=\frac1{Q(D)}.}
 \label{rm:eq:confluence-ew-master}
\end{equation}

The first equality publishes the same coupling in two gauge
orientations; the second recovers it as a sum of compliances. The determinant
\(D\) governs the even mixture, and \(\rho\) preserves the odd orientation.

The Higgs coordinate incorporates the latter:

\begin{equation}
 \lambda_H=\frac{Q(D)}{8(1-D)}D^{-3}\rho^{5/3}.
 \label{rm:eq:confluence-higgs}
\end{equation}

Defining

\begin{equation}
 \Xi=\frac{8(1-D)D^3\lambda_H}{Q(D)},
 \label{rm:eq:confluence-xi}
\end{equation}

is reconstructed

\begin{equation}
 \rho=\Xi^{3/5},
 \qquad
 \widehat Z=
 \frac{F(\sqrt D\,\Xi^{3/10})}
      {F(\sqrt D\,\Xi^{-3/10})}.
 \label{rm:eq:confluence-vacuum-higgs}
\end{equation}

The equation \cref{rm:eq:confluence-vacuum-higgs} provides a cross-check between scalar orientation and vacuum impedance. Its force comes from the parallel construction of both readers on \(X\).

\subsection{2nd Confluence: Action, Gravity, and Electrogravitational Form}

Let

\begin{equation}
 Q_P=E_P,
 \qquad
 \ell_P,
 \qquad
 Q_P\ell_P=\hbar c,
 \qquad
 \frac{\ell_P}{Q_P}=\frac{G}{c^4}.
 \label{rm:eq:confluence-planck-pair}
\end{equation}

The product is the exchange quantum \(\hbar c\); the quotient is the
gravitational compliance. We introduce

\begin{equation}
 \mathscr T_P=\frac{Q_P}{\ell_P}=\frac{c^4}{G},
 \qquad
 \mathscr C_P=\frac{\ell_P}{Q_P}=\frac{G}{c^4}.
 \label{rm:eq:confluence-tension-compliance}
\end{equation}

The Einstein equation assumes the constitutive form

\begin{equation}
 G_{\mu\nu}+\Lambda g_{\mu\nu}
 =8\pi\mathscr C_P T_{\mu\nu}.
 \label{rm:eq:confluence-einstein-constitutive}
\end{equation}

For a charge \(q\), the electric coordinate lifted to the energy plane
is

\begin{equation}
 \mathcal Q(q)=Q_P\sqrt{\frac{Z}{4\pi\hbar}}\,q.
 \label{rm:eq:confluence-electric-coordinate}
\end{equation}

Newton and Coulomb are then brought together in the bilinear form

\begin{equation}
 \boxed{
 V(r)=\frac{\mathscr C_P}{r}
 \bigl(\mathcal Q_1\mathcal Q_2-E_1E_2\bigr).}
 \label{rm:eq:confluence-newton-coulomb}
\end{equation}

The cone \(\mathcal Q^2=E^2\) reproduces the static condition of
extremality in the Einstein--Maxwell chart.  The sheet product of the vacuum
governs the causal cone; its quotient governs the electric lift and
orientation.

The cross-identity with the electroweak sector is

\begin{equation}
 \boxed{
 \frac{\hbar^2}{8\pi\ell_P^2}
 \frac{(\kappa/\mu)}{e^2}
 =\frac1{4\pi\alpha_{\HMT}}
 =\frac1{g^2}+\frac1{g'^2}.}
 \label{rm:eq:confluence-einstein-maxwell-ew}
\end{equation}

This equality preserves its value under inversion of the action sheet:
\(G\) and \(e^2\) transform contragrediently, whereas
\(G\hbar\), \(Ge^2\), and \(\ell_P^2\) remain invariant.

\subsection{3rd confluence: geometry, information and mixing}

Let \(N=N_{90}+N_{120}\) be the total number of boundary excitations and

\begin{equation}
 D_{90/120}=90N_{90}+120N_{120}.
 \label{rm:eq:confluence-modular-provenance}
\end{equation}

The area reads \(N\), whereas the modular holographic Hamiltonian reads
\(D_{90/120}\). Since

\begin{equation}
 \det\begin{pmatrix}1&1\\90&120\end{pmatrix}=30,
 \label{rm:eq:confluence-modular-determinant}
\end{equation}

the pair reconstructs exactly

\begin{equation}
 N_{90}=4N-\frac{D_{90/120}}{30},
 \qquad
 N_{120}=\frac{D_{90/120}}{30}-3N.
 \label{rm:eq:confluence-sector-reconstruction}
\end{equation}

The total geometry and the weighted modular energy therefore preserve
quantity and provenance.

The angular double-layer state

\begin{equation}
 \rho_y=\frac{e^{-yR}}{2\cosh y}
 \label{rm:eq:confluence-rhoy}
\end{equation}

has even entropy and oriented divergence

\begin{equation}
 S(\rho_y)=\log(2\cosh y)-y\tanh y,
 \qquad
 D(\rho_y\Vert\rho_{-y})=2y\tanh y.
 \label{rm:eq:confluence-info-orientation}
\end{equation}

The involution \(y\mapsto-y\) is represented in the flavour sector by

\begin{equation}
 R_{\rm CKM}
 =M_{\rm CKM}\operatorname{diag}(1,-1,1)M_{\rm CKM}^{-1},
 \qquad
 R_{\rm CKM}^2=I,
 \qquad
 \det R_{\rm CKM}=-1.
 \label{rm:eq:confluence-rckm}
\end{equation}

Barbero reads the odd component of the transition \(6\to8\), entropy
reads its even distribution, and \(R_{\rm CKM}\) realizes the same inversion in the
mixing codomain.  They are coordinated representations of a common hinge,
with their own domains and observables.

\subsection{4th Confluence: Screen, Torsion, and Helical Memory}

In the quotient lorentziano \(Q_{\square}\), the soldadura and the respuis
positiva satisfy

\begin{equation}
 \boxed{
 (\beta_m^{\square})^*\Lambda_m^{\square}\beta_m^{\square}
 =\frac{56}{81}I_{Q_{\square}}.}
 \label{rm:eq:confluence-screen-energy}
\end{equation}

In the deep register, the screw operator checks

\begin{equation}
 \boxed{
 V^{-1}\mathscr SV=(2+\sqrt3)\ii\,\mathscr S.}
 \label{rm:eq:confluence-screw}
\end{equation}

The first identity measures the minimum response cost of the screen; the
second transports orientation and depth. Catalan appears as the
second moment of the quaternary component. If
\(N_{\rm dep}\delta_n=n\delta_n\) denotes the depth operator of the
helical record, then

\begin{equation}
 G_{\rm Cat}=\Tr(N_{\rm dep}^{-2}\mathcal Q_{\chi,-4}).
 \label{rm:eq:confluence-catalan}
\end{equation}

Perron's theorem

\begin{equation}
 M_n=\left(\frac{a_n}{a_0}\right)^{-6}
 \label{rm:eq:confluence-sixth-law}
\end{equation}

provides the exact covariance of a quadratic spin density.  The
Cartan realizer must transport
\cref{rm:eq:confluence-screen-energy,rm:eq:confluence-screw} to a current and a
contorsion whose homogeneous mode satisfies
\cref{rm:eq:confluence-sixth-law}.  This chain places the physical interface in a unique and falsifiable arrow.

\subsection{Confluence between clock periodicity, horizon energy, and information cost}

The clock \(108\) fixes

\begin{equation}
 E_P=k_B\Theta_{\rm clk}\log3.
 \label{rm:eq:confluence-planck-trit}
\end{equation}

For the de Sitter patch of radius \(R=\ell_P\), the vacuum energy,
temperature, and entropy satisfy

\begin{equation}
 E_\Lambda=T_H S_H=\frac{E_P}{2}
 =\frac{k_B\Theta_{\rm clk}\log3}{2}.
 \label{rm:eq:confluence-horizon-half}
\end{equation}

From here follow

\begin{equation}
 E_\Lambda t_P=\frac{\hbar}{2},
 \qquad
 E_\Lambda T_{108}=\frac{h}{2},
 \qquad
 \frac{S_H}{k_B}=\pi.
 \label{rm:eq:confluence-half-action}
\end{equation}

The vacuum volume, the horizon area, the cycle temperature, and the information cost of a trit publish the same energy. The multiplicity \(e^{S_H/k_B}=e^\pi\) coincides exactly with the expansive eigenvalue of \(e^{\pi J_\varphi}\); promotion of this equality to an equivalence of microstates requires an intertwiner that preserves measure and dynamics.

\subsection{Fibrewise confluence of HMT–MD responses: vacuum, action, gravity, information, mixing, memory, and cosmology}

\begin{theorem}[Confluence of HMT--MD responses]
Given the enriched state \(X\) and the primitives declared in each reader, the publications of vacuum, action, gravity, information, mixing, screen, helical memory, and cosmology are the projections or compositions defined on the fiber product of graphs \cref{rm:eq:master-fiber-product}. Their cross-relations are generated by \cref{rm:eq:confluence-ew-master,rm:eq:confluence-newton-coulomb,%
rm:eq:confluence-einstein-maxwell-ew,rm:eq:confluence-sector-reconstruction,%
rm:eq:confluence-info-orientation,rm:eq:confluence-screen-energy,%
rm:eq:confluence-screw,rm:eq:confluence-horizon-half}.
Every numerical evaluation is terminal, and every
oriented memory remains in the factor that generates it.
\end{theorem}

\begin{proof}
The construction of \(\mathfrak M_X\) fixes the common antecedent. The
vacuum equations follow from the functional calculation of \(T\); the
electroweak equations follow from \(D\) and \(Q(D)\); the electrogravitational form follows
from the two identities of the Planck pair; the reconstruction \(90/120\)
results from inverting an integral matrix of determinant \(30\); the informational orientation
follows from the logarithm of \(\rho_y\); the screen energy
results from the unique response in \(Q_{\square}\); the screw follows from the
two Weyl and Pell covariances; and the horizon equality follows from the
geometric formulas of de Sitter specialized at \(R=\ell_P\). No
proof uses the terminal decimal to select the branch. The horizontal
independence theorem prevents cross-relations from being converted into
unproved factorizations.
\end{proof}

\subsection{Corollary of metrological confluence between the HMT and SI realizations}

Every output admits the inseparable quintuple

\begin{equation}
 \mathcal K(P)=
 \bigl(
 \text{antecedent},
 \text{equation},
 \text{invariant},
 \text{meaning},
 \text{terminal value}
 \bigr).
 \label{rm:eq:confluence-quintuple}
\end{equation}

The volume's metrological atlas enumerates these quintuples. Its
probative content is as follows: the antecedent fixes the genealogy; the equation
fixes the operation; the invariant demonstrates what survives; the meaning
makes the physical reading explicit; the terminal value enables external
recognition. Separating one column destroys causal information; preserving all
five permits insertion of the result into a chart of units without converting
that chart into a foundation.

\begin{proofstatusbox}
The fiber product, the algebraic identities, and the horizontal-independence theorem constitute an exact formalization. The internal equations cited retain the proved status of the chapters containing their complete proofs. The Einstein--Maxwell, de Sitter, Einstein--Cartan, and electroweak readings remain at their declared interfaces when they require an additional realizer. The theorem organizes those layers without promoting them by numerical similarity.
\end{proofstatusbox}

\begin{falsifierbox}
Confluence is refuted if two components attributed to the same element of \(\mathfrak M_X\) actually arise from incompatible states, if one of the generating identities fails, if a terminal evaluation enters into the selection of its own antecedent, or if a memory distinguished by another reader is lost. A purported horizontal factorization is refuted by any pair of states with the same intermediate reader and different final outputs. Physical promotions are additionally tested against their own realizers, gauge identities, Bianchi identities, refinement, and boundary conditions.
\end{falsifierbox}
